% ==============================================================================
% reports/chapters/03_du_lieu_va_bai_toan.tex
% CHƯƠNG 3: TẬP DỮ LIỆU VÀ ĐỊNH NGHĨA BÀI TOÁN
% ==============================================================================

\chapter{Tập Dữ Liệu và Định Nghĩa Bài Toán}
\label{chap:du_lieu}

\section{Tổng Quan Tập Dữ Liệu CICIDS2017 Tuesday}

Tập dữ liệu CICIDS2017 \cite{sharafaldin2018generating,cicids2017dataset} được xây dựng bởi Viện An ninh mạng Canada (CIC), Đại học New Brunswick (UNB). Đây là một trong những bộ dữ liệu chuẩn mực và được sử dụng rộng rãi để nghiên cứu và đánh giá các hệ thống phát hiện xâm nhập mạng.

Tập dữ liệu ghi lại lưu lượng mạng trong 5 ngày làm việc liên tiếp (từ thứ Hai 03/07/2017 đến thứ Sáu 07/07/2017) trên hạ tầng mạng mô phỏng gồm mạng nạn nhân và mạng tấn công bên ngoài. Trong đó, ngày thứ Ba (Tuesday, 04/07/2017) được dành riêng để thực hiện các kịch bản tấn công dò quét mật khẩu (Brute Force):
\begin{itemize}
    \item \textbf{Tấn công buổi sáng}: \FTP\ thực hiện dò quét mật khẩu dịch vụ FTP (Cổng 21) từ 09:20 đến 10:20.
    \item \textbf{Tấn công buổi chiều}: \SSH\ thực hiện dò quét mật khẩu dịch vụ SSH (Cổng 22) từ 14:00 đến 15:00.
\end{itemize}

\section{Nguồn Gốc Dữ Liệu và Xác Minh Tính Toàn Vẹn SHA-256}

Dữ liệu thô được sử dụng trong bài tập lớn là tệp CSV trích xuất chính thức từ gói \texttt{GeneratedLabelledFlows.zip} của UNB:
\begin{itemize}
    \item \textbf{Đường dẫn}: \texttt{data/raw/Tuesday-WorkingHours.pcap\_ISCX.csv}
    \item \textbf{Mã băm SHA-256 đã xác minh}:
    \begin{center}
        \small\texttt{\DataShaFull}
    \end{center}
    \item \textbf{Tổng số dòng thô ban đầu}: \DataRawRows\ bản ghi flow (tương ứng 445,910 dòng trên đĩa tính cả 1 dòng tiêu đề).
\end{itemize}

Việc lưu trữ và kiểm tra mã băm SHA-256 cho phép xác minh tệp đầu vào có cùng nội dung byte với tệp đã ghi nhận mã băm, hỗ trợ tính tái lập thực nghiệm trên cùng tệp dữ liệu nguồn.

\section{Phục Hồi Dòng Thời Gian 24 Giờ (Timestamp Reconstruction)}

Trong quá trình khám phá dữ liệu (EDA), chúng tôi nhận thấy trong tệp Tuesday được sử dụng, các timestamp quan sát được không đủ thông tin để phân biệt một số giờ buổi sáng/buổi chiều theo dòng thời gian kỳ vọng:
\begin{itemize}
    \item Định dạng chuỗi thời gian thô được lưu dưới dạng \texttt{dd/MM/yyyy} kết hợp với \texttt{h:mm} hoặc \texttt{hh:mm}.
    \item Các giá trị giờ thô quan sát được chỉ gồm hai nhóm:
    \begin{equation}
        \mathcal{H}_{\text{raw}} = \{ 1, 2, 3, 4, 5, 8, 9, 10, 11, 12 \}
    \end{equation}
    mà không có hậu tố chỉ định AM/PM.
    \item Các bản ghi diễn ra vào buổi chiều (từ 13:00 đến 17:00) xuất hiện dưới dạng 1:00 đến 5:00.
\end{itemize}

Để phục hồi đúng thứ tự thời gian phục vụ phép phân chia theo thời gian (Time-based Split), dự án áp dụng quy tắc tái dựng thời gian \texttt{reconstruct\_timestamp()} tại mã nguồn \path{scripts/prepare_splits.py}:
\begin{equation}
    H_{\text{24h}} = \begin{cases}
        H_{\text{raw}} + 12 & \text{nếu } H_{\text{raw}} \in \{1, 2, 3, 4, 5\} \\
        H_{\text{raw}} & \text{nếu } H_{\text{raw}} \in \{8, 9, 10, 11, 12\}
    \end{cases}
\end{equation}

\noindent\textbf{Nguyên tắc kiểm soát rò rỉ dữ liệu (No-Leakage Policy):} Đây là một quy tắc tái dựng của pipeline dự án, được thiết lập hoàn toàn độc lập với nhãn tấn công (\texttt{Label}) và cổng đích (\texttt{Destination Port}), chỉ dựa trên phân tích dòng thời gian của toàn bộ traffic trong ngày Tuesday (thời gian làm việc từ 8:30 sáng đến 17:00 chiều).

\section{Định Nghĩa Bài Toán và Không Gian Nhãn}

Bài tập lớn định nghĩa bài toán phân loại là \textbf{phân loại nhị phân} nhằm nhận diện lưu lượng bất thường:
\begin{equation}
    y = \begin{cases}
        0 & \text{nếu flow là lưu lượng bình thường (\BENIGN)} \\
        1 & \text{nếu flow là tấn công Brute Force (\FTP\ hoặc \SSH)}
    \end{cases}
\end{equation}

Chi tiết ánh xạ nhãn được trình bày trong Bảng \ref{tab:label_mapping}.

% ==============================================================================
% reports/tables/manual/tab_label_mapping.tex
% ==============================================================================
\begin{table}[htbp]
\centering
\caption{Quy tắc ánh xạ nhãn bài toán phân loại nhị phân phát hiện tấn công dò quét mật khẩu.}
\label{tab:label_mapping}
\adjustbox{max width=\textwidth}{%
\begin{tabular}{lclp{7.5cm}}
\toprule
\textbf{Nhãn gốc trong dữ liệu} & \textbf{Nhãn $y$} & \textbf{Xử lý bài toán} & \textbf{Mô tả bản chất an ninh mạng} \\
\midrule
\BENIGN & 0 & Lớp thường (\Normal) & Lưu lượng mạng của các tác vụ và giao dịch hợp lệ của người dùng nội bộ. \\
\FTP & 1 & Lớp tấn công (\Attack) & Hành vi dò quét thử mật khẩu liên tục đối với dịch vụ FTP (Cổng đích 21). \\
\SSH & 1 & Lớp tấn công (\Attack) & Hành vi dò quét thử mật khẩu tự động đối với dịch vụ SSH (Cổng đích 22). \\
\bottomrule
\end{tabular}%
}
\end{table}


Mặc dù bài toán phân loại chính là nhị phân, thông tin chi tiết về dạng tấn công (\texttt{Subtype}) vẫn được bảo toàn trong một cột siêu dữ liệu riêng biệt để theo dõi tỷ lệ phát hiện chi tiết từng dạng tấn công (Subtype Recall).

\section{Quy Trình Làm Sạch Dữ Liệu Thô}

Tệp dữ liệu thô trải qua các bước làm sạch kỹ thuật:
\begin{enumerate}
    \item \textbf{Loại bỏ flow có thời lượng không hợp lệ}: Loại bỏ \DataInvalidDurationRows\ bản ghi có $\text{Flow Duration} \le 0$ hoặc bị thiếu (NaN).
    \item \textbf{Loại bỏ bản ghi trùng lặp}: Loại bỏ \DataDuplicatesRows\ dòng bản ghi trùng lặp hoàn toàn trên tất cả các cột dữ liệu.
    \item \textbf{Loại bỏ cột đặc trưng trùng lặp}: Cột \texttt{Fwd Header Length.1} có giá trị trùng lặp 100\% với cột \texttt{Fwd Header Length} được loại bỏ khỏi không gian đặc trưng.
    \item \textbf{Xử lý giá trị vô cùng}: Có \DataInfinityValues\ giá trị vô cùng ($\pm\infty$, chủ yếu phát sinh do phép chia cho thời lượng bằng 0 trong các cột \texttt{Flow Bytes/s} và \texttt{Flow Packets/s}) được chuyển đổi thành giá trị \texttt{NaN} để xử lý qua bộ điền giá trị thiếu (Imputer).
\end{enumerate}

Sau khi làm sạch kỹ thuật, tổng số bản ghi hợp lệ đưa vào quy trình phân tách là \textbf{\DataCleanRows\ dòng}.

\section{Dòng Thời Gian Tấn Công Thực Tế và Phân Bổ Mẫu}

Bảng \ref{tab:dataset_splits} tổng hợp chi tiết số lượng mẫu và phân bổ nhãn qua các tập sau khi áp dụng chiến lược phân tách theo thời gian (được trình bày chi tiết ở Chương 4).

\begin{table}[htbp]
\centering
\caption{Thống kê phân bổ dữ liệu các tập Train, Validation và Test theo chiến lược Temporal Split v1.}
\label{tab:dataset_splits}
\adjustbox{max width=\textwidth}{%
\begin{tabular}{lrrrrrr}
\toprule
\textbf{Tập dữ liệu} & \textbf{Tổng số flow} & \textbf{BENIGN} & \textbf{Tấn công} & \textbf{FTP-Patator} & \textbf{SSH-Patator} & \textbf{Tỷ lệ Attack (\%)} \\
\midrule
Train & 90,273 & 85,139 & 5,134 & 5,134 & 0 & 5.69\% \\
Validation & 203,826 & 199,522 & 4,304 & 2,418 & 1,886 & 2.11\% \\
Test & 144,157 & 140,425 & 3,732 & 0 & 3,732 & 2.59\% \\
\midrule
Purge/Embargo & 7,632 & 6,967 & 665 & 386 & 279 & 8.71\% \\
\textbf{Tổng số hợp lệ} & \textbf{445,888} & \textbf{432,053} & \textbf{13,835} & \textbf{7,938} & \textbf{5,897} & \textbf{3.10\%} \\
\bottomrule
\end{tabular}%
}
\end{table}


Phân tích dòng thời gian cho thấy đặc điểm phân bố tự nhiên của dữ liệu CICIDS2017 Tuesday:
\begin{itemize}
    \item \textbf{Tập Train} ($e < 10:00$): Chứa \DataTrainBenign\ mẫu bình thường và \DataTrainFTP\ mẫu \FTP. \textbf{Không có mẫu \SSH\ nào (0 mẫu)}.
    \item \textbf{Tập Validation} ($10:02 \le s$ và $e < 14:30$): Chứa đoạn kết của chiến dịch FTP (\DataValidationFTP\ mẫu) và giai đoạn đầu của chiến dịch SSH (\DataValidationSSH\ mẫu).
    \item \textbf{Tập Test} ($s \ge 14:32$): Chứa \DataTestBenign\ mẫu bình thường và \DataTestSSH\ mẫu \SSH. \textbf{Không có mẫu \FTP\ nào (0 mẫu)}.
\end{itemize}

Đặc tính phân bố này tạo ra bài toán đánh giá khả năng \textbf{tổng quát hóa sang biến thể mới (kịch bản FTP $\to$ SSH)}. Trong thiết lập này, mô hình hoàn toàn không được huấn luyện trên mẫu \SSH\ nào trong tập Train; tập Validation chứa 1,886 mẫu SSH tham gia vào quá trình lựa chọn siêu tham số, trước khi mô hình được đánh giá độc lập trên tập Test chỉ chứa biến thể SSH.
